\subsection*{把复利的小步改成切向小步}
普通复利每步按当前量增加一小部分。若把增量乘上 $\mathrm i$，新增部分不再沿原方向，而是垂直于原方向。一步变成乘 $1+\mathrm i h$，不断重复，就得到一条绕原点前进的折线轨迹。

每步到底是不是旋转？它的模长为 $\sqrt{1+h^2}>1$，所以还夹带一点向外扩张。图像若只强调绕圈，容易隐藏这个有限步误差。精确的圆周运动来自步长趋零后的极限，而不是任意有限步都在圆上。

\subsection*{为什么误差在连续极限消失}
固定总时间 $t$，分成 $n$ 步，模长倍率是 $(1+t^2/n^2)^{n/2}$。单步误差是二阶的，总步数只是一阶增加；取对数后约为 $t^2/(2n)$，趋向零，所以总模长趋向一。

角度每步增加 $\arctan(t/n)$，总角度为 $n\arctan(t/n)$，趋向 $t$。模长和角度两个条件合起来，解释了极限为何是 $\mathrm e^{\mathrm it}$。这里用到的局部近似必须在小步长下理解。

这个例子也说明欧拉数值法可能产生什么误差：真实系统守恒半径，有限切向步却慢慢漂离。把步数加倍，观察偏差变化，会比单看一条漂亮轨迹更有教学价值。Lean 只核验单步模长恒等式，收敛论证仍需另外完成。
